\section{总结与延伸}

本讲可以压缩成一句话：\textbf{训练 helpful chatbot 不是选择一个 optimizer，而是共同设计 data distribution、preference protocol、training objective 与 measurement system。} H4 human-data experiments 与 Zephyr synthetic-data pipeline 展示两条可行路线；GPT-4 evaluator audit 则证明，若 measurement 与 training data 同源，排行榜会奖励 style matching 而非真实目标。

\subsection{十一条可执行结论}

\begingroup
\footnotesize
\setlength{\columnsep}{1.2em}
\begin{multicols}{2}
\begin{enumerate}[itemsep=0.05em,topsep=0.1em,parsep=0pt,leftmargin=*]
  \item 报告 dataset provenance：谁写 prompt、谁写 completion、谁过滤。
  \item 用 task/length/turn distribution 描述数据，不只报 example count。
  \item 区分 demonstration、preference pair 与 policy rollout 的计数单位。
  \item SFT 先建立 instruction-following；DPO/RLHF 通常是 refinement，不是替代地基。
  \item Human-written 与 synthetic data 的错误相关性不同，应保留异构 holdout。
  \item Reward model 必须在新 generators、难分 pairs 与 held-out tasks 上评估。
  \item 同时报告 task metrics、human preference、LLM judge 与 response statistics。
  \item 遇到 metric reversal 不要挑赢家，要定位 evaluator preference。
  \item LLM judge 必须做 order swap、length control、human calibration 与 task stratification。
  \item Red teaming 测尾部风险，不能被平均 helpfulness score 代替。
  \item 成本、规模和 leaderboard 都必须标注日期、单位与 protocol。
\end{enumerate}
\end{multicols}
\endgroup

\begin{importantbox}{最终系统视角}
一条 alignment pipeline 可以写成闭环：
\[
\text{Data source}\rightarrow\text{Training objective}\rightarrow\text{Candidate behavior}
\rightarrow\text{Evaluator}\rightarrow\text{Next data round}.
\]
Evaluator 的偏差会决定下一轮收什么数据，因此 measurement error 会被重新写回 training distribution。真正的工程目标不是一次赢榜，而是让闭环可审计、可校准、可纠错。
\end{importantbox}

\subsection{拓展阅读}

\begingroup
\scriptsize
\begin{itemize}[itemsep=0pt,topsep=0.15em,parsep=0pt]
  \item Ouyang et al., \href{https://arxiv.org/abs/2203.02155}{Training language models to follow instructions with human feedback}：SFT、reward modeling 与 PPO/RLHF 基础流程。
  \item Wang et al., \href{https://arxiv.org/abs/2212.10560}{Self-Instruct}；Ding et al., \href{https://arxiv.org/abs/2305.14233}{UltraChat}；Li et al., \href{https://arxiv.org/abs/2303.17760}{CAMEL}：三种 synthetic instruction/dialogue generation 路线。
  \item Tunstall et al., \href{https://arxiv.org/abs/2310.16944}{Zephyr: Direct Distillation of LM Alignment}：dSFT、AI feedback 与 dDPO。
  \item Zheng et al., \href{https://arxiv.org/abs/2306.05685}{Judging LLM-as-a-Judge with MT-Bench and Chatbot Arena}：LLM judge、position bias 与 human agreement。
  \item Hugging Face H4, \href{https://huggingface.co/blog/llm-leaderboard}{LLM evaluator bias analysis} 与 \href{https://huggingface.co/blog/red-teaming}{red-teaming workflow}：课堂 slides 的官方实践入口。
\end{itemize}
\endgroup

\end{document}
